% =====================================================================
\chapter{Vectores en el plano y en el espacio}
\label{tema:1}
% =====================================================================

Un vector es, a la vez, una lista de números y una flecha con longitud y dirección. Este tema relaciona las dos caras: las operaciones se definen con coordenadas y se interpretan geométricamente. El producto escalar permite medir longitudes y ángulos; el producto vectorial, áreas y volúmenes. La presentación sigue las sesiones~1 a~7 de la unidad~1 del curso 18.02SC del MIT~\cite{mit1802}, las secciones 2.1--2.4 de OpenStax~\cite{openstax3} y la sección~11.1 de Strang~\cite{strangcalc}.

\section{Vectores y operaciones}

\begin{definicion}[Vectores de $\R^n$]\label{def:vector}
Para $n=2$ o $n=3$, un \emph{vector} de $\R^n$ es una lista ordenada $\vu=(u_1,\dots,u_n)$ de números reales, sus \emph{componentes}. Para $\vu,\vv\in\R^n$ y $\lambda\in\R$ se definen la \emph{suma} y el \emph{producto por un escalar} componente a componente:
\[
\vu+\vv=(u_1+v_1,\dots,u_n+v_n),\qquad \lambda\vu=(\lambda u_1,\dots,\lambda u_n).
\]
El vector nulo es $\vcero=(0,\dots,0)$ y $-\vu=(-1)\vu$. En $\R^3$ se escribe $\vi=(1,0,0)$, $\vj=(0,1,0)$, $\vk=(0,0,1)$, de modo que $(u_1,u_2,u_3)=u_1\vi+u_2\vj+u_3\vk$.
\end{definicion}

Las propiedades de la suma y del producto por escalares (asociativa, conmutativa, distributivas, neutro y opuesto) se heredan de las de los números reales, porque se opera componente a componente. Un \emph{punto} y un \emph{vector} se escriben igual; la diferencia es de interpretación. Si $P$ y $Q$ son puntos, el vector $\overrightarrow{PQ}=Q-P$ es la flecha que va de $P$ a $Q$; dos flechas con la misma longitud, dirección y sentido representan el mismo vector, aunque partan de puntos distintos.

Geométricamente, $\vu+\vv$ es la diagonal del paralelogramo de lados $\vu$ y $\vv$ (o, equivalentemente, la flecha que se obtiene colocando $\vv$ a continuación de $\vu$), y $\lambda\vu$ es una flecha paralela a $\vu$, $\abs{\lambda}$ veces más larga y con el sentido invertido si $\lambda<0$.

\addtocounter{definicion}{1}%

\section{Producto escalar, norma y ángulo}

\begin{definicion}[Producto escalar y norma]\label{def:escalar}
El \emph{producto escalar} de $\vu,\vv\in\R^n$ es el número
\[
\vu\cdot\vv=u_1v_1+\dots+u_nv_n=\sum_{i=1}^{n}u_iv_i .
\]
La \emph{norma} (o longitud) de $\vu$ es $\norm{\vu}=\sqrt{\vu\cdot\vu}=\sqrt{u_1^2+\dots+u_n^2}$, y la \emph{distancia} entre $\vu$ y $\vv$ es $\norm{\vu-\vv}$. Un vector de norma~1 es \emph{unitario}; si $\vu\neq\vcero$, el vector $\vu/\norm{\vu}$ es el unitario en la dirección de $\vu$.
\end{definicion}

En $\R^2$ y $\R^3$ la norma es la longitud de la flecha, por el teorema de Pitágoras (aplicado dos veces en $\R^3$).

\begin{proposicion}[Propiedades del producto escalar]\label{prop:escalar}
Para todos $\vu,\vv,\vw\in\R^n$ y $\lambda\in\R$:
\begin{enumerate}
  \item simetría: $\vu\cdot\vv=\vv\cdot\vu$;
  \item linealidad: $(\vu+\vw)\cdot\vv=\vu\cdot\vv+\vw\cdot\vv$ y $(\lambda\vu)\cdot\vv=\lambda(\vu\cdot\vv)$;
  \item positividad: $\vu\cdot\vu\ge 0$, con igualdad si y solo si $\vu=\vcero$;
  \item $\norm{\lambda\vu}=\abs{\lambda}\,\norm{\vu}$.
\end{enumerate}
\end{proposicion}

Desarrollando con (i) y (ii) se obtiene la identidad, que se usará muchas veces,
\begin{equation}\label{eq:t1-desarrollo}
\norm{\vu\pm\vv}^2=(\vu\pm\vv)\cdot(\vu\pm\vv)=\norm{\vu}^2\pm2\,\vu\cdot\vv+\norm{\vv}^2 .
\end{equation}

\begin{teorema}[Desigualdad de Cauchy--Schwarz]\label{teo:cs}
Para todos $\vu,\vv\in\R^n$, $\abs{\vu\cdot\vv}\le\norm{\vu}\,\norm{\vv}$. La igualdad se da si y solo si uno de los vectores es múltiplo del otro.
\end{teorema}

\begin{corolario}[Desigualdad triangular]\label{cor:triangular}
Para todos $\vu,\vv\in\R^n$, $\norm{\vu+\vv}\le\norm{\vu}+\norm{\vv}$.
\end{corolario}

Por Cauchy--Schwarz, si $\vu,\vv\neq\vcero$ el cociente $\vu\cdot\vv/(\norm{\vu}\norm{\vv})$ está en $[-1,1]$, lo que justifica la definición siguiente.

\begin{definicion}[Ángulo y ortogonalidad]\label{def:angulo}
El \emph{ángulo} entre dos vectores no nulos $\vu$ y $\vv$ es el único $\theta\in[0,\pi]$ tal que
\[
\cos\theta=\frac{\vu\cdot\vv}{\norm{\vu}\,\norm{\vv}} .
\]
Dos vectores son \emph{ortogonales} (o perpendiculares), $\vu\perp\vv$, si $\vu\cdot\vv=0$.
\end{definicion}

\begin{proposicion}[El ángulo algebraico es el geométrico]\label{prop:coseno}
Si $\vu,\vv\neq\vcero$ forman, como flechas con origen común, un ángulo geométrico $\theta$, entonces $\vu\cdot\vv=\norm{\vu}\norm{\vv}\cos\theta$.
\end{proposicion}

Así, el signo de $\vu\cdot\vv$ informa del ángulo: positivo si es agudo, nulo si es recto y negativo si es obtuso. Si $\vu\perp\vv$, la identidad~\eqref{eq:t1-desarrollo} se reduce al \emph{teorema de Pitágoras} $\norm{\vu+\vv}^2=\norm{\vu}^2+\norm{\vv}^2$.

\addtocounter{definicion}{1}%

\section{Proyecciones}

\begin{definicion}[Proyección ortogonal]\label{def:proyeccion}
Sean $\vu,\vv\in\R^n$ con $\vv\neq\vcero$. La \emph{componente escalar} de $\vu$ en la dirección de $\vv$ es $\operatorname{comp}_{\vv}\vu=\vu\cdot\vv/\norm{\vv}$, y la \emph{proyección ortogonal} de $\vu$ sobre $\vv$ es el vector
\[
\proj_{\vv}\vu=\frac{\vu\cdot\vv}{\vv\cdot\vv}\,\vv .
\]
\end{definicion}

\begin{proposicion}[Descomposición ortogonal]\label{prop:descomposicion}
Sea $\vv\neq\vcero$. Todo $\vu$ se escribe de forma única como $\vu=\vu_\parallel+\vu_\perp$ con $\vu_\parallel$ múltiplo de $\vv$ y $\vu_\perp\perp\vv$; además, $\vu_\parallel=\proj_{\vv}\vu$. Entre todos los múltiplos $t\vv$, el más cercano a $\vu$ es $\proj_{\vv}\vu$.
\end{proposicion}

\addtocounter{figure}{1}%

\addtocounter{definicion}{1}%

\section{Producto vectorial y producto mixto}

\begin{definicion}[Producto vectorial]\label{def:vectorial}
El \emph{producto vectorial} de $\vu,\vv\in\R^3$ es el vector
\[
\vu\times\vv=(u_2v_3-u_3v_2,\ u_3v_1-u_1v_3,\ u_1v_2-u_2v_1)
=\begin{vmatrix}\vi&\vj&\vk\\u_1&u_2&u_3\\v_1&v_2&v_3\end{vmatrix},
\]
donde el determinante es simbólico y se desarrolla por la primera fila (\cref{def:det3}).
\end{definicion}

\begin{proposicion}[Propiedades del producto vectorial]\label{prop:vectorial}
Para todos $\vu,\vv,\vw\in\R^3$ y $\lambda\in\R$:
\begin{enumerate}
  \item anticonmutatividad: $\vv\times\vu=-\vu\times\vv$; en particular, $\vu\times\vu=\vcero$;
  \item linealidad: $(\vu+\vw)\times\vv=\vu\times\vv+\vw\times\vv$ y $(\lambda\vu)\times\vv=\lambda(\vu\times\vv)$;
  \item ortogonalidad: $(\vu\times\vv)\cdot\vu=(\vu\times\vv)\cdot\vv=0$;
  \item $\norm{\vu\times\vv}=\norm{\vu}\norm{\vv}\sin\theta$, donde $\theta$ es el ángulo entre $\vu$ y $\vv$; por tanto, $\norm{\vu\times\vv}$ es el área del paralelogramo de lados $\vu$ y $\vv$;
  \item $\vi\times\vj=\vk$, $\vj\times\vk=\vi$, $\vk\times\vi=\vj$.
\end{enumerate}
\end{proposicion}

La dirección de $\vu\times\vv$ es, pues, perpendicular al plano que contienen $\vu$ y $\vv$, y su sentido viene dado por la \emph{regla de la mano derecha}: $(\vu,\vv,\vu\times\vv)$ está orientado como $(\vi,\vj,\vk)$. El producto vectorial no es asociativo: $(\vi\times\vi)\times\vj=\vcero$, mientras que $\vi\times(\vi\times\vj)=\vi\times\vk=-\vj$.

\begin{definicion}[Producto mixto]\label{def:mixto}
El \emph{producto mixto} de $\vu,\vv,\vw\in\R^3$ es el número $\vu\cdot(\vv\times\vw)$.
\end{definicion}

\begin{proposicion}[Volumen]\label{prop:volumen}
$\abs{\vu\cdot(\vv\times\vw)}$ es el volumen del paralelepípedo de aristas $\vu$, $\vv$ y $\vw$ con un vértice común. En particular, el producto mixto es nulo si y solo si los tres vectores están en un mismo plano por el origen.
\end{proposicion}

\addtocounter{definicion}{1}%

\addtocounter{definicion}{1}%

\addtocounter{definicion}{1}%

